\documentclass[11pt]{article}
% Shared preamble for all agentic-payments memos.
% Usage:
%   \documentclass[11pt]{article}
%   % Shared preamble for all agentic-payments memos.
% Usage:
%   \documentclass[11pt]{article}
%   % Shared preamble for all agentic-payments memos.
% Usage:
%   \documentclass[11pt]{article}
%   \input{_preamble}
%   \title{...} \author{...} \date{...}
%   \begin{document} \maketitle ... \end{document}

\usepackage[margin=1in]{geometry}
\usepackage[T1]{fontenc}
\usepackage[utf8]{inputenc}
\usepackage{lmodern}
\usepackage{microtype}
\usepackage{booktabs}
\usepackage{array}
\usepackage{longtable}
\usepackage{enumitem}
\usepackage{xcolor}
\usepackage{hyperref}
\usepackage{amsmath}
\usepackage{amssymb}
\usepackage{graphicx}
\usepackage{caption}
\usepackage{subcaption}
\usepackage{listings}

\hypersetup{
  colorlinks=true,
  linkcolor=blue!55!black,
  citecolor=blue!55!black,
  urlcolor=blue!55!black
}

\setlist[itemize]{topsep=3pt,itemsep=2pt,leftmargin=1.4em}
\setlist[enumerate]{topsep=3pt,itemsep=2pt,leftmargin=1.7em}

\definecolor{codebg}{RGB}{246,247,242}
\lstset{
  basicstyle=\ttfamily\footnotesize,
  backgroundcolor=\color{codebg},
  frame=single,
  rulecolor=\color{black!12},
  breaklines=true,
  showstringspaces=false,
  columns=fullflexible,
  keepspaces=true,
}

\newcommand{\code}[1]{\texttt{\detokenize{#1}}}
\newcommand{\risk}{\operatorname{risk}}

% Experimental result callout (used by data-driven memos).
\newenvironment{result}
  {\par\medskip\noindent\textbf{Result.}\quad\itshape}
  {\upshape\par\medskip}
\newenvironment{hypothesis}
  {\par\medskip\noindent\textbf{Hypothesis.}\quad}
  {\par\medskip}
\newenvironment{experiment}
  {\par\medskip\noindent\textbf{Experiment.}\quad}
  {\par\medskip}

%   \title{...} \author{...} \date{...}
%   \begin{document} \maketitle ... \end{document}

\usepackage[margin=1in]{geometry}
\usepackage[T1]{fontenc}
\usepackage[utf8]{inputenc}
\usepackage{lmodern}
\usepackage{microtype}
\usepackage{booktabs}
\usepackage{array}
\usepackage{longtable}
\usepackage{enumitem}
\usepackage{xcolor}
\usepackage{hyperref}
\usepackage{amsmath}
\usepackage{amssymb}
\usepackage{graphicx}
\usepackage{caption}
\usepackage{subcaption}
\usepackage{listings}

\hypersetup{
  colorlinks=true,
  linkcolor=blue!55!black,
  citecolor=blue!55!black,
  urlcolor=blue!55!black
}

\setlist[itemize]{topsep=3pt,itemsep=2pt,leftmargin=1.4em}
\setlist[enumerate]{topsep=3pt,itemsep=2pt,leftmargin=1.7em}

\definecolor{codebg}{RGB}{246,247,242}
\lstset{
  basicstyle=\ttfamily\footnotesize,
  backgroundcolor=\color{codebg},
  frame=single,
  rulecolor=\color{black!12},
  breaklines=true,
  showstringspaces=false,
  columns=fullflexible,
  keepspaces=true,
}

\newcommand{\code}[1]{\texttt{\detokenize{#1}}}
\newcommand{\risk}{\operatorname{risk}}

% Experimental result callout (used by data-driven memos).
\newenvironment{result}
  {\par\medskip\noindent\textbf{Result.}\quad\itshape}
  {\upshape\par\medskip}
\newenvironment{hypothesis}
  {\par\medskip\noindent\textbf{Hypothesis.}\quad}
  {\par\medskip}
\newenvironment{experiment}
  {\par\medskip\noindent\textbf{Experiment.}\quad}
  {\par\medskip}

%   \title{...} \author{...} \date{...}
%   \begin{document} \maketitle ... \end{document}

\usepackage[margin=1in]{geometry}
\usepackage[T1]{fontenc}
\usepackage[utf8]{inputenc}
\usepackage{lmodern}
\usepackage{microtype}
\usepackage{booktabs}
\usepackage{array}
\usepackage{longtable}
\usepackage{enumitem}
\usepackage{xcolor}
\usepackage{hyperref}
\usepackage{amsmath}
\usepackage{amssymb}
\usepackage{graphicx}
\usepackage{caption}
\usepackage{subcaption}
\usepackage{listings}

\hypersetup{
  colorlinks=true,
  linkcolor=blue!55!black,
  citecolor=blue!55!black,
  urlcolor=blue!55!black
}

\setlist[itemize]{topsep=3pt,itemsep=2pt,leftmargin=1.4em}
\setlist[enumerate]{topsep=3pt,itemsep=2pt,leftmargin=1.7em}

\definecolor{codebg}{RGB}{246,247,242}
\lstset{
  basicstyle=\ttfamily\footnotesize,
  backgroundcolor=\color{codebg},
  frame=single,
  rulecolor=\color{black!12},
  breaklines=true,
  showstringspaces=false,
  columns=fullflexible,
  keepspaces=true,
}

\newcommand{\code}[1]{\texttt{\detokenize{#1}}}
\newcommand{\risk}{\operatorname{risk}}

% Experimental result callout (used by data-driven memos).
\newenvironment{result}
  {\par\medskip\noindent\textbf{Result.}\quad\itshape}
  {\upshape\par\medskip}
\newenvironment{hypothesis}
  {\par\medskip\noindent\textbf{Hypothesis.}\quad}
  {\par\medskip}
\newenvironment{experiment}
  {\par\medskip\noindent\textbf{Experiment.}\quad}
  {\par\medskip}


\title{\textbf{Pitch framing}\\
\large Memo 01}
\author{Bloomsbury Tech -- Agentic Payments Hackathon Build}
\date{April 25, 2026}

\begin{document}
\maketitle

\textbf{Date:} 2026-04-25

\section{The two-minute pitch (working draft)}
\begin{quote}
Visa, Mastercard, and Stripe Radar built fraud detection on a billion \emph{human} transactions. They detect anomalies relative to a human behavioural baseline --- geography, time-of-day, value bands. AI agents transact a thousand times more often, optimise gas to the wei, have no spatial pattern, and behave non-stationarily as the underlying models update. Existing fraud models are blind to them.  We pulled 100,000 recent Base L2 transactions and built behavioural embeddings --- here's the cluster plot. Agents and humans separate cleanly. Within the agent cluster, our anomaly detector flags compromised wallets the existing systems miss, and our interpretability layer gives a one-sentence reason for every flag.  We're Bloomsbury Tech --- a quantitative ML lab that publishes on causal inference and graph ML. \textbf{This is a statistical problem, not an LLM problem.} Stripe Radar can't ship this because they make money on agent transactions; Visa is on a 24-month release cycle; Sardine, Persona, Alloy were built for human-only data. We have an 18-month structural lead, and every month of agent traffic compounds our model.  We start as the fraud signal layer for one mid-market PSP or acquirer. We expand into agent reputation scoring, AML for stablecoin flows, and dispute evidence generation. The eventual buyer is Stripe, Visa, Adyen, or Sardine --- but by then we own the data.
\end{quote}
\section{Four beats}
\begin{enumerate}
  \item \textbf{Precise gap.} Existing fraud models trained on human distributions; agents are structurally different.
  \item \textbf{Demo evidence.} UMAP cluster plot + anomaly explanations on real(ish) Base data.
  \item \textbf{Moat.} Statistical not LLM; data compounds; incumbents structurally slow.
  \item \textbf{Expansion.} Fraud \textbackslash{}\$\textbackslash{}to\textbackslash{}\$ reputation \textbackslash{}\$\textbackslash{}to\textbackslash{}\$ AML \textbackslash{}\$\textbackslash{}to\textbackslash{}\$ disputes \textbackslash{}\$\textbackslash{}to\textbackslash{}\$ buyer is the network.
\end{enumerate}
\section{Blomfield rubric (what he scores)}
\begin{itemize}
  \item Precise problem statement (no vision waffle)
  \item An \emph{unusual} insight (the human-baseline fraud-model blindness)
  \item Evidence of velocity (working demo in 8 hours)
  \item Defensible moat (not LLM; data compounds; channel conflict for incumbents)
  \item Founder credibility (ICML-grade ML lab, not a vibe-coded MVP)
\end{itemize}
\section{Open framing decision}
"Fraud detection" = easier pitch, smaller TAM. "Behavioural risk intelligence for agent traffic" = same product, bigger headroom (fraud, reputation, AML, disputes). \textbf{Default: pitch the broader frame, demo the fraud slice.} Blomfield watched Monzo go prepaid-card \textbackslash{}\$\textbackslash{}to\textbackslash{}\$ full bank; he'll track the expansion arc.

\end{document}
